\chapter{反抗分析}
《剧场版魔法少女小圆【新篇】叛逆的物语》是一个颇为棘手的分析对象。
与动画一样，这部电影的符号密度极高（也就是说，片中的画面和事件往往能够引发多重解读）。
不过，动画通常以冷峻、空旷，几乎不带生气的空间环境作为背景，并借此凸显魔女结界中繁复而异质的视觉表达；
而《叛逆》则几乎在整个叙事过程中都维持着高度复杂的视觉组织。
换句话说，剧场版的每一幅画面都像动画中的画面一样，承载着丰富的解读可能，而且全片包含的意象之多，远远超出相当于五集半的片长给人的预期！

因此，分析这部电影很容易让人不知从何下手。
若想在一章的篇幅内将其作为一个完整的艺术整体加以把握，几乎是不现实的。
要真正较为充分地讨论它，通常只能采取两种路径：
其一，是选定一个核心主题，沿着这一主题考察它在影片中的发展与变化；
但这样做势必会牺牲其他主题的讨论，除非它们与该主题有直接关联。
其二，是集中分析某一个具体场景，尽可能展开其内部的复杂结构；
但这种方法同样会使其他场景退居幕后，只有当它们影响到所选场景的理解时，才会被纳入讨论。

我们可以从一个看似简单的问题切入：
所谓“叛逆”，究竟指向什么？
它是否指丘比对小圆所创造的新世界秩序的反抗？
是否指焰在尚未意识到自己正是结界创造者时，对结界中那个世界的抗拒？
又或者，它指的是焰对小圆本人的反叛？
进一步说，影片本身是否也可以被理解为一次叛逆？
如果可以，那么它所反叛的对象究竟是谁或是什么？

即使是“片名究竟意味着什么？”
这样表面上十分简单的问题，也并不会带来一个明确的答案，反而会不断引出新的疑问。
而这些疑问中的每一个，都可能得到多种有充分依据的解释；
其中任何一种解释，都足以单独展开成一篇完整的论述。
因此，本书接下来的章节将由一组围绕《叛逆》展开的专题文章构成。
有些文章会聚焦于具体场景的细读，有些则会追索某一主题的展开，或分析某个角色的变化过程。
但无论是哪种方式，这些讨论最终都指向同一个核心问题：所谓“叛逆”，究竟是在反抗什么？
在我看来，这种从不同角度逐步逼近问题的方式，是理解《叛逆的物语》唯一较为可行的路径。
因为这部电影并不轻易服从单一、封闭的阐释框架，它本身就是一部不断逃逸、不断挑战分析的作品。

而反抗的近义词是……你应该明白了\footnote{%
    译注：原文为``And a near-synonym for defiance is... well, you get the idea.''
         译者认为这是作者的搞怪，``defiance（反抗）''与``rebellion（叛逆）''同义，故做此译。
}。

我们不妨从一个具体而短小的场景入手，来说明这部电影为何如此难以阐释。
这个场景就是片尾彩蛋。
在此之前，片尾本身以高度风格化的方式重现了影片的情节，焰和小圆分别位于缓缓滚动的职员表两侧，仿佛被隔开般。
但到了最后，两人却牵起彼此的手，一同奔向远方。
相较于正片结尾处焰近乎挑衅、强作姿态的表现，这一画面显得意外地温柔，甚至带有某种希望色彩。
这一幕是在预示未来的发展，还是只是焰的梦境？
她们最终消失在远方，究竟意味着两人终将摆脱当前的困境，还是意味着她们真正相伴的可能性正在消逝？

不过，这一点并不重要，因为片尾彩蛋本身已经对前面的图像构成了否定。
（或者它真的否定了吗？如果片尾呈现的是焰的梦境，而彩蛋是现实，或者片尾是预示，而彩蛋则是她的恐惧……）
这里有必要先考察一下彩蛋通常具有的作用。
彩蛋多见于大型商业动作系列电影中，漫威电影宇宙尤其将这种做法发展得极为成熟；它也常见于喜剧片。
它的功能通常有两种：
第一类是收束影片前文埋下的笑点，作出最后回应。
例如，《空前绝后满天飞！》中那位出租车乘客，就是这类用法的典型例子\cite{airplane-amazon-video}。
第二类则是为系列的后续作品制造期待，同时埋下关于下一部情节走向的线索。
比如塞缪尔·杰克逊\footnote{%
    塞缪尔·杰克逊（1948年12月21日—），美国男演员及监制。
}出场表示想和你谈谈“复仇者联盟计划”\cite{iron-man-amazon-video}。

然而，这里的彩蛋似乎既不属于前者，也不属于后者。
它的主要功能似乎是提出新的疑问，并动摇观众对焰在字幕前最后几个场景中自我呈现方式的理解。
诚然，从形式上看，这一段也可以被理解为对《魔法少女小圆》剧场版系列中某一反复母题的“收束”，即“草地中的两把椅子”。
在前两部总集篇电影的片头中，小圆与焰并肩坐在花草地中的白色椅子上，彼此依偎，姿态亲密而轻快。
在《叛逆》彩蛋中，焰独自坐在相似的椅子上，位置却从开阔的花草地变成了悬崖边缘。
椅子、断崖与半月共同构成一种被切割的视觉结构，仿佛整幅画面被一分为二：
月亮失去了另一半，世界也失去了另一半；
而那被切除的部分，正包括小圆以及原本属于她的那把椅子。
原本属于小圆的部分，只剩下一片空洞的黑暗。

这一画面究竟意味着什么？
是焰已经下定决心，甚至要把小圆也置于自己的对立面吗？
是她内心悔意的显现吗？
还是说，它只是给观众一次残酷提醒，让我们意识到那些已经失去的东西？

在影片结尾，焰看起来几乎已经掌握了绝对的控制权。
她拥有听命于自己的使魔军队，也能够改写沙耶香的记忆，切断她与人鱼魔女形态的联系；
她甚至可以阻断小圆与她的“佛性”，也就是“圆环之理”之间的联系。
此前，焰已经通过改写现实创造了这个新的世界，因此，把她视为这一宇宙内部最强大的存在，并非没有根据。
不过，她与“圆环之理”之间究竟孰强孰弱，仍然难以断定。
因为影片中有两个对象是她似乎无法彻底掌控的：一个是她自己，另一个则是小圆。
片尾字幕前，小圆一度开始想起自己作为“圆环之理”的“佛性”本质，请注意焰随即显露出的惊慌。

那么，在彩蛋中，为什么焰看到丘比靠近时会显得惊讶呢？
她脸上的表情可以解读为不安或希望；
考虑到椅子的象征，她是否一时以为那是小圆？
她希望是小圆，还是害怕是小圆？
她怎么会不知道那是丘比呢？

再来看看丘比的状态：
它的毛发凌乱，身体颤抖。
无论是在动画还是在这部影片中，丘比眼睛的极近距离特写都反复出现。
这类镜头通常用来提醒观众：丘比始终在观察、监视和评估。
因此，它的眼部特写往往强化了其不祥感和威胁性。
但画面中的丘比毛发打结，眼神黯淡，瞳孔发黑并微微颤动。
它第一次显露出某种情绪反应，然而这种变化只会进一步加深观众的不安感。
因为在这一刻，即便是向来冷静、空洞、难以撼动的丘比，也被扭曲成了某种“异常”的存在。
它不再是那个令人畏惧的操控者，而成了一个可怜、狼狈的对象，被焰改写宇宙时展现的力量所击垮。
对它和它的族群来说，这是最坏的结果：
小圆虽然最终对丘比产生了恐惧和不信任，但她并不真正具备强烈憎恨他人的能力。
焰则截然不同。
她的内心充满愤怒与悲伤，因此，如果她将这种情绪发泄到孵化者这一族群整体之上，尤其是发泄到见泷原这个丘比身上，也完全符合她此时的性格与心理状态。

然而，从片尾彩蛋的叙事逻辑来看，丘比眼部的极端特写应该带有某种伏笔意味。
在整段彩蛋中，这一镜头最接近通常意义上的片尾彩蛋：
它仿佛恐怖电影结尾的经典处理——观众以为杀手已经死去，镜头却突然显示他睁开了眼睛，暗示威胁并未真正消失。
问题在于，丘比依旧像过去一样难以读解。
它究竟是在暗中筹划对焰的反击，还是只是在冷静观察、等待时机？
又或者，它真的已经被焰彻底击溃了？
它此刻近乎狼狈而可怜的形象，是否正表明焰已经剥夺了它作为反派的叙事位置？

再说那段舞蹈。
焰与她新的灵魂宝石共舞，其舞蹈风格和配乐都让人联想起电影开篇附近，她那段芭蕾式的变身场景。
这颗宝石的外形酷似国际象棋中的“王”：虽是棋盘上几乎毫无攻击力的棋子，却决定了整盘棋的生死。
这是否意味着，身为“后”（棋盘上最强棋子）的焰，正在与她世界中最重要的人物——如今已形同虚设的小圆——翩翩共舞？
这种解读是可以得到支持的：
这颗灵魂宝石由焰的旧灵魂宝石碎片和一轴与小圆发色相同的线制成（这一场景出现在此前），但前提是，我们将那根线视为小圆本人，或者她与焰之间纽带的象征。
另一种同样合理的解读是：
这根线代表“具身化、作为个体的小圆”与“无处不在、无形无相、作为抽象概念的小圆”之间的联结。
若取此解，焰便并非在思念她的“另一半”，而是沉浸于囚禁小圆的病态满足之中。
换言之，这两种解读的分歧在于：焰究竟是在困惑与痛苦之上佯装坚强，还是已沦为一名令人不寒而栗、控制欲极端的跟踪狂。

在舞蹈的尾声，焰侧身坠下悬崖。
她坠落的身姿让人想起焰变成魔女时，小圆同样以类似方式，从椅子上侧身倾倒的画面；
这一呼应显然指向小圆的自我牺牲，以及焰因未能阻止她而愈发深重的悔恨。
那么，焰是否在通过复现小圆的举动，寻求借此与小圆相聚？
是在牺牲自己以确保小圆不再需要这么做吗？
还是在嘲弄小圆的牺牲，借此表明自己已经堕落到如此地步，以至于连最初推动她走向堕落的爱本身也失去了意义？
也可能，这只是一个绝望至极的角色所作出的无望姿态：
一种指向自杀，却注定无法真正完成的举动？
焰死后宇宙会变成什么样，影片并没有说明。
但几乎可以确定的是，至少小圆会重新与自己的“佛性”本质——
即圆环之理——
相连，而焰一直试图阻断的正是这种联系。
这一画面也可能与上述解读并存，或是取代上述解读指向塔罗牌中的“愚者”。
在许多“愚者”牌中，通常描绘着一名愚人正要迈下悬崖，身边伴着一只白色小动物\footnote{%
    可参见颇具影响力的韦特塔罗牌中的“愚者”图像。
}\cite{rider-waite-fool-image}。

我们固然可以详尽探讨上述问题，以及其他相关问题（例如，月相恰好处于沙耶香魔女化时的几近新月与焰魔女化时的满月之间的正中位置，这一事实究竟有何寓意）。% TODO figure
这样的分析虽需耗费数千言，且最终除确认该场景具有刻意为之的模糊性外，难有定论，但确实可行。
然而，问题关键不在此处。
关键在于，这一幕仅占两小时全片中的九十秒，甚至并非视觉或符号密度最高的一段（窃以为，最高密度的段落当属焰化为魔女直至她与小圆共同击碎孵化者结界之间的若干片段）。

不，重点是：
这部电影的符号密度极高，并且充满歧义，这对解读工作提出了极大的挑战。

很好。
那我们就来尝试解读吧。
